% HY562 / CS-562 Fall 2026, course project report (6 pages, LNCS format, references not counted).
% Use the Springer LNCS template: on Overleaf, "Springer LNCS template", or llncs.cls from
% https://www.springer.com/gp/computer-science/lncs/conference-proceedings-guidelines
\documentclass[runningheads]{llncs}
\usepackage{graphicx,booktabs,amsmath,url}

\begin{document}
\title{Project A1: ShapeStream with Sketches}   % your project, team and variant
\author{Name One \and Name Two \and Name Three \and Name Four \and Name Five}
\institute{Computer Science Department, University of Crete}
\maketitle

\begin{abstract}
Four to six sentences: the problem, what you built, the main number, what it means.
\end{abstract}

\section{Introduction}
The problem in one paragraph; your variant and why it is interesting; your contributions as a short list.

\section{Problem and background}
Definitions you use (from the project description); the lecture material you build on; related work from the reading list.

\section{The common baseline}
Project A: design of the exact incremental engine (state, update rules for the three hard cases, Spark operators).
Project B: the pipeline and the two baselines as you ran them; anything you had to fix.

\section{Your variant}
Design, parameters, complexity or error guarantees where they exist. One worked example
(Project A: one epoch on a few entities; Project B: one question from summary to query to answer).

\section{Experiments}
Set-up (machine, Spark settings, data, models). The minimum grid of the project description.
Every number in this section must be a row of the shared results table.

\section{Discussion}
What worked, what did not and why; threats to validity; what you would do next.

\section*{Use of AI assistants}
Which tools, for what. (Course policy, Lecture 1.)

\bibliographystyle{splncs04}
\bibliography{refs}
\end{document}
